\chapterhead{3}{Drops on a Cold Glass}{Finding where the water came from}\label{ch3}

Fill a glass with ice water on a warm day and wait. Little drops appear on the \emph{outside} of the glass. They grow, join, and run down. Where did that water come from?

There are two easy ideas. The water might have \textbf{leaked through the glass} from the inside. Or it might have come \textbf{out of the air} around the glass. Both could be true. A test can tell them apart.

\begin{center}
\begin{tikzpicture}
  \draw[line width=1.4pt,color=ink] (0,0) -- (0,2.6) (1.6,0) -- (1.6,2.6) (0,0) -- (1.6,0);
  \fill[sky!45] (0.03,0.03) rectangle (1.57,1.9);
  \foreach \x/\y in {-0.18/0.8,-0.2/1.4,1.78/0.6,1.8/1.3,1.78/1.9,-0.17/2.0}{\fill[sky!80!black] (\x,\y) circle (0.07);}
  \draw[-{Stealth},line width=1.4pt,color=sky!80!black] (2.9,1.5) -- (1.9,1.2);
  \draw[-{Stealth},line width=1.4pt,color=sky!80!black] (-1.3,1.5) -- (-0.3,1.2);
  \node[font=\sffamily\scriptsize,anchor=west] at (3.0,1.5) {vapour from the air?};
  \node[font=\sffamily\scriptsize,anchor=east] at (-1.4,1.5) {leak from inside?};
\end{tikzpicture}
\end{center}

\begin{tryit}{Cold glass and colour}
\textbf{You need:} two glasses that are the same, cold water, ice, food colouring or a drop of ink, a clock, paper towel.
\begin{enumerate}[leftmargin=1.6em]
\item Dry the outside of both glasses. Fill one with ice water and add a few drops of food colouring. Fill the other with room-temperature water and the same colouring.
\item Set them side by side. Count the seconds until you first see drops on the outside. Stop after 10 minutes if there are none.
\item Dab a drop with a white paper towel. Is it coloured or clear?
\item Do it three times.
\end{enumerate}
\par\medskip
\renewcommand{\arraystretch}{1.6}
\begin{tabular}{|l|c|c|c|c|}
\hline
 & \textbf{Try 1} & \textbf{Try 2} & \textbf{Try 3} & \textbf{Average}\\ \hline
\textbf{Cold glass: seconds to first drops} & & & & \\ \hline
\textbf{Room glass: seconds to first drops} & & & & \\ \hline
\textbf{Colour of drops (clear/coloured)} & & & & \\ \hline
\end{tabular}
\end{tryit}
\adult{Food colouring stains clothes and hands. Use cold water and ice only.}

\begin{worked}{Practice numbers, not from a real test}
Cold glass: 90, 80, 100 seconds. Total 270, average $270\div3=90$ seconds, spread 20.\\
Room glass: no drops in 10 minutes, all three tries.\\
The drops from the cold glass were clear every time.
\end{worked}

\subhead{What the test says}
If the glass leaked, the drops would be coloured like the water inside. They were clear. So the drops did not come from the inside. The air was the likely source. The room glass fits that too, since without a cold surface no drops formed.

Air always holds some invisible water vapour. When air touches something cold, the air next to it cools. Cool air can hold less vapour, so some vapour turns back into tiny drops of liquid water. That change, vapour to liquid, is \textbf{condensation}. It is evaporation in reverse.

\clearpage

\begin{center}
\begin{tikzpicture}
  \node[bigcir,minimum size=2.4cm] (a) at (0,0) {}; \node[font=\sffamily\scriptsize,align=center] at (0,0) {air near\\the glass};
  \node[bigcir,minimum size=2.4cm] (b) at (5.4,0) {}; \node[font=\sffamily\scriptsize,align=center] at (5.4,0) {drops on\\the glass};
  \draw[moveA] (a) -- node[above,note] {condensation} (b);
  \node[note,anchor=west,text width=3.1cm,align=left] at (7.1,0) {\textbf{kept:} the water moves from air to glass\\\textbf{set aside:} how many drops, and their sizes};
\end{tikzpicture}
\end{center}

\begin{multix}[hmm]\multixbot[hmm]\tcblower
The drops are on the outside of the glass and the water is on the inside, so the glass must be leaking a tiny bit. I am a careful robot. Did I miss a test?
\end{multix}

\trust{We repeated the cold glass three times and added colour to the water inside to see whether the drops carried it.}
{The drops were clear every time, which is hard to explain if they came through the glass.}
{How much water came from the air, or whether a tiny, invisible leak also happened. The colour test only rules out a leak big enough to carry visible colour.}

\begin{keepbox}
We said the drops came from the air. \textbf{We kept} \wline[3.4cm] \textbf{and we set aside} \wline[3.4cm].
\end{keepbox}

\subhead{Fog on a spoon}
\begin{tryit}{Breath on a cold spoon}
Put one metal spoon in the fridge for 10 minutes and leave another at room temperature. Breathe gently on each. Which one fogs? \wline[5cm]\\[4pt]
Does the fog stay? What happens as the spoon warms? \wline[5cm]
\end{tryit}

\begin{worked}{Practice result}
The cold spoon fogged at once and cleared in about 10 seconds as it warmed. The room spoon barely fogged. Your breath carries extra vapour, and the cold spoon cools it enough for drops to form.
\end{worked}

\subhead{The cling film again}
Look back at the covered cup in Chapter~1. Drops on the inside of the film are condensation too. In a covered cup the vapour has nowhere to go, so it meets the film and turns back to drops. Some run back into the water.

\clearpage

\begin{puzzle}{Average the times}
Cold glass: 90, 80 and 100 seconds. Total \wline[1.6cm] Average \wline[1.6cm]~s
\end{puzzle}

\begin{puzzle}{Colour clue}
Blue water inside a cold glass gave clear drops outside. Which idea does that support? Circle one. \\
\textbf{leaked through} \quad \textbf{came from the air} \quad \textbf{cannot tell}
\end{puzzle}

\begin{puzzle}{No drops}
Why did the room-temperature glass stay dry outside? \wline[6cm]
\end{puzzle}

\begin{puzzle}{Unfair test}
Mia put a cold glass in the kitchen and a warm glass outside in the sun. Only the cold glass had drops. Name two things that were different. \wline[6cm]
\end{puzzle}

\begin{puzzle}{Cold things}
Name two cold things that collect drops. \wline[5cm]
\end{puzzle}

\clearpage

\begin{puzzle}{Move the water}
Suppose a cold glass collected 2 spoonfuls of water from a room. Fill in the circles.
\begin{center}
\begin{tikzpicture}
  \node[cir] (a) at (0,0) {50}; \node[cir] (b) at (5,0) {0};
  \draw[moveA] (a) -- node[above,note] {2 spoonfuls} (b);
  \node[note] at (0,-0.8) {air vapour before}; \node[note] at (5,-0.8) {glass drops before};
  \node[cir,dashed] at (0,-2.2) {}; \node[cir,dashed] at (5,-2.2) {};
  \node[note] at (0,-3.0) {air after}; \node[note] at (5,-3.0) {glass after};
\end{tikzpicture}
\end{center}
\end{puzzle}

\begin{puzzle}{A sentence for the data}
Mark S if the data supports it, G if it is a guess about why, and B if it is too big for the data.\\
(a) ``The drops from the cold glass were clear.'' \wline[1cm]\\
(b) ``The drops came out of the air.'' \wline[1cm]\\
(c) ``All drops everywhere come from air.'' \wline[1cm]
\end{puzzle}

\begin{puzzle}{Morning window}
On a cold morning the inside of a window has drops. Is the glass cold or warm? \wline[3cm] Where could the water come from? \wline[4cm]
\end{puzzle}

\begin{puzzle}{Why the film?}
In the covered cup, why might drops form on the cling film? \wline[6cm]
\end{puzzle}
\clearpage
